% ---------------------------------------------------------------------------
% howtoread.tex
% ---------------------------------------------------------------------------
\chapter*{How to Read This Booklet}
\addcontentsline{toc}{chapter}{How to Read This Booklet}
\markboth{How to Read This Booklet}{How to Read This Booklet}

\section*{Typographic conventions}

\begin{itemize}
  \item \code{Monospaced text} is code: an identifier, a keyword, a file name
        or a command.
  \item \kw{Coloured monospaced text} marks a language keyword under
        discussion.
  \item \file{Green monospaced text} is a file in the accompanying source
        tree.
  \item \tool{Sans-serif} names are tools: \tool{QuestaSim}, \tool{Verilator},
        \tool{cocotb}, \tool{GTKWave}.
\end{itemize}

\section*{The three kinds of box}

\begin{gotcha}{What a trap looks like}
A trap marks a mistake that VHDL experience makes you \emph{more} likely to
commit, not less.
\end{gotcha}

\begin{tipbox}{What a recipe looks like}
A recipe is a pattern to copy. When the text says ``write it like this'',
this is the box it is in.
\end{tipbox}

\begin{notebox}{What a note looks like}
A note supplies background: where a construct came from, what the standard
actually says, why a tool behaves the way it does. Skip these on a first
reading and come back when something surprises you.
\end{notebox}

\begin{ruleofthumb}
A rule of thumb is one line you should be able to recite from memory.
\end{ruleofthumb}

\section*{If you have been told to write a testbench by Friday}

\Cref{ch:tb} is self-contained. Read it, take the layered structure, and
ignore \cref{ch:uvm} entirely until somebody asks you for reuse across
projects. If your simulator is \tool{Verilator} or you have no licence at
all, or your golden model is already written in Python, see
\refsimulationbook once the structure here makes sense.

\section*{If you already know the design-side language}

Start at \cref{ch:toolkit}: it is the vocabulary (classes, dynamic types,
constrained randomisation, functional coverage, assertions, dynamic
threads) that the rest of this booklet assumes. None of it is
synthesisable, and none of it has a VHDL counterpart, so there is nothing to
translate here, only to learn.

\section*{The companion booklets}

This is the \textbf{Verification Cookbook}. \refdesignbook covers the
design-side language this booklet assumes you already know.
\refsimulationbook covers \tool{Verilator}, \tool{QuestaSim} and
Python/\tool{cocotb} testbenches. \refexamplebook works one CORDIC design
through all of it, including the testbenches built here.
